\documentclass[11pt,a4paper]{article}
\usepackage[utf8]{inputenc}
\usepackage[T1]{fontenc}
\usepackage{soul}
\newcommand\dunderline[3][-2pt]{{%
  \sbox0{#3}%
  \ooalign{\copy0\cr\rule[\dimexpr#1-#2\relax]{\wd0}{#2}}}}
\usepackage[left=2cm, right=2cm, top=2cm, bottom=2cm]{geometry}
\usepackage{enumitem}
\usepackage{hyperref}
\hypersetup{
	colorlinks=true,
	linkcolor=cyan,
	filecolor=blue,
	urlcolor=blue,
	citecolor=blue,
}
\usepackage{xcolor}
\usepackage{parskip}

\pagestyle{empty}

\begin{document}

\begin{center}
    {\LARGE \dunderline{1.5pt}{\textbf{Curriculum Vitae}}}
\end{center}

\vspace{0.5em}

\dunderline{1.5pt}{\large{\textbf{Name:} Yang-yang Tan}}

\vspace{1cm}

\dunderline{1.5pt}{\large{\textbf{A. Education}}}
\vspace{0.5cm}

\noindent
\textbf{Dalian University of Technology}, Dalian, Liaoning Province, China\\
Ph.D., Theoretical Physics\\
Thesis: \textit{Towards Real-Time Dynamics of Critical Phenomena in Strongly Correlated Systems}\\
% Advisor: Prof.\ Wei-jie Fu\\
Date of Graduation: July 2025

\vspace{0.6em}

\noindent
\textbf{Dalian University of Technology}, Dalian, Liaoning Province, China\\
B.S., Physics\\
% Thesis: \textit{The Deep Learning Method in Gluon Spectra Function}\\
% Advisor: Prof.\ Wei-jie Fu\\
Date of Graduation: July 2019

\vspace{1cm}

\dunderline{1.5pt}{\large{\textbf{B. Research/Employment History}}}
\vspace{0.5cm}

\noindent
\begin{tabular}{@{}p{3.8cm}p{12cm}@{}}
Apr.\ 2026 -- Present & Postdoctoral Researcher, The University of Tokyo, Tokyo, Japan \\[0.3em]
Dec.\ 2025 -- Mar.\ 2026 & Visiting Scholar, Tsinghua University, Beijing, China \\[0.3em]
% Sep.\ 2019 -- Jul.\ 2025 & Ph.D.\ Student, Dalian University of Technology, Dalian, China \\
\end{tabular}

\vspace{1cm}

\dunderline{1.5pt}{\large{\textbf{C. Fellowships / Awards}}}
\vspace{0.5cm}

\vspace{1cm}

\dunderline{1.5pt}{\large{\textbf{D. Other Activities}}}
\vspace{0.5cm}

\large{\textbf{(1). Academic Collaboration}}
\normalsize\vspace{1ex}

\begin{itemize}[leftmargin=2em]
\item \textbf{DM-LFT Collaboration} (Diffusion Model-Lattice Field Theory): G.\ Aarts, S.\ Chen, D.\ E.\ Habibi, A.\ Ipp, R.\ Kapust, B.\ Lucini, D.\ M\"uller, J.\ M.\ Pawlowski, T.\ R.\ Ranner, Y.-y.\ Tan, L.\ Wang, W.\ Wang, K.\ Zhou, and Q.\ Zhu. (2026)

\item \href{https://fqcd-collaboration.github.io}{\textbf{fQCD Collaboration}} (Functional QCD): J.\ Braun, Y.-r.\ Chen, W.-j.\ Fu, F.\ Gao, C.\ Huang, F.\ Ihssen, K.\ Kockler, Y.\ Lu, J.\ M.\ Pawlowski, F.\ Rennecke, F.\ Sattler, J.\ Stoll, Y.-y.\ Tan, Z.-n.\ Wang, R.\ Wen, J.\ Wessely, S.\ Yin, H.-w.\ Zheng and N.\ Zorbach. (2026)
\end{itemize}

\vspace{0.4cm}

\large{\textbf{(2). Technical Skills}}
\normalsize\vspace{1ex}

\noindent GitHub Profile: \href{https://github.com/Yangyang-Tan}{\textcolor{blue}{https://github.com/Yangyang-Tan}}

\begin{itemize}[leftmargin=2em]
\item Wolfram Language \& Mathematica --- Expert, 12 years of experience
\item Julia --- Proficient, 5 years of experience
\item Python (PyTorch) --- Proficient, 2 years of experience
\item Large-scale parallel computing and GPU computing (based on Julia) --- Proficient, 5 years of experience
\end{itemize}

\end{document}
